\scanpage{148}
Faisant $M=N=N'$, on trouve qu'une loi de composition bilinéaire sur $V(M)$ correspond à un hom
\[
\Delta:M\mathbin{\text{---}}M\otimes M
\]
(appelé aussi ``application diagonale''). \deleted{On dit que}\inserted{on dit que $M$ \unsure{devt|devient} une coalgèbre.}
$\Delta$ est \unsure{coassociative|associative}, \inserted{resp.} commutative, \inserted{resp.} unitaire si la loi de composition sur $V(M)$ l'est.
\deleted{\unsure{l'unité}}\inserted{Cela} \unsure{s'explicite} par les diagrammes bien connus. Pour l'unitarité, on note que $e\in V(M)(k)$ s'identifie à une \uline{coaugmentation} ou counité
\[
M\xrightarrow{\unsure{\eta}}k
\]
(satisfaisant des conditions qu'on n'écrit pas, qui la définissent de façon unique).

Par contre, une \deleted{\unsure{la}} coaugmentation (multiplicative)
$V(M)\mathbin{\text{---}}O_k\inserted{=V(k)}$ correspond à un hom $k\mathbin{\text{---}}M$, ou encore un $\mathrm{él}^t$ de $M$ (appelé unité)
la condition ici est
\[
\Delta e=e\otimes e.
\]
Notons que par suite de la formule
\[
V\left(\lim_i M_i\right)=\deleted{\unsure{\prod_i}}\varprojlim_i V(M)
\]
(pour cat. d'indices quelc.), une lim. quelc. de coalgèbres est une coalgèbre, et les \deleted{structures}\inserted{axiomes} d'ass, d'unité, de comm.
\margin{147}
